\chapter{Developer guide and limitations}\label{ch:c9}\label{ch:dev}
This chapter is for those who maintain or extend the code.

\begin{plain}
This chapter is for programmers: which file does what, how one time step passes through the code, how to add a new
exchange or grid type, and how to test changes. Users of the coupling can skip it.
\end{plain}

\section{Source files}
\begin{table}[!htbp]\centering
\caption{Source files.}
\begin{tabular}{P{0.36\linewidth}P{0.588\linewidth}}\toprule\hdr
File & Content\\\midrule
\cmd{GroundwaterModel.h/.cpp} & class \cmd{CGroundwaterModel}: parse-time setters, model builder, per-step exchange, outputs, hotstart\\
\cmd{GWExternal.cpp} & coupling to an existing MODFLOW 6 model: reads the simulation through the MF6 API, copies it, rewrites the time discretization (and renumbers periods on restarts), links HRUs (cell outlines with \cmd{:MF6CRS}, polygon file, or overlap weights), writes \cmd{RCH\_RAVEN}/\cmd{DRN\_RAVEN}, maps the stream package to reaches, collects the other packages' flows\\
\cmd{GWUtil.h} & small helpers shared by the coupling's files (absolute paths, MODFLOW array writer, number formatting, the strict number rule of the \cmd{.rvg} and \cmd{.rvc})\\
\cmd{GWGrid.h/.cpp} & \cmd{CGWGrid}: cells of one layer as polygons with areas, centres, neighbours (face length, centre-to-face distances, face direction); generators for regular (rotation, variable spacing), quadtree (with its size estimated before it is built), Voronoi and imported grids; point location, candidate search, exact clipping of polygons and segments; cache serialization\\
\cmd{GWGeometry.h/.cpp} & \cmd{CGWProjection} (equal-area projection), \cmd{CGWCRS} (projections of existing models: Transverse Mercator, Albers, Lambert conformal conic), \cmd{CGWRaster} (ESRI ASCII and GeoTIFF, bilinear sampling), minimal JSON/GeoJSON reader, the EPSG code classification, polygon clipping and areas\\
\cmd{GWGeoTIFF.h/.cpp} & GeoTIFF reader without external libraries: classic TIFF and BigTIFF, strips and tiles, LZW, DEFLATE (own inflate) and PackBits decoders, predictors, georeferencing and GeoTIFF keys; every offset and count checked against the file\\
\cmd{GWEPSGTable.h} & the geographic (longitude/latitude) EPSG codes, generated from the EPSG registry by \cmd{tools/dev/make\_epsg\_geographic.py} of the extra repository\\
\cmd{MF6Engine.h/.cpp} & \cmd{CMF6Engine}: loads \cmd{libmf6} at run time (\cmd{dlopen}/\cmd{LoadLibrary}); initialize, time step, Newton loop, memory access; switches to the MODFLOW folder around every call\\
\cmd{ParseGWFile.cpp} & \cmd{.rvg} parser\\\bottomrule
\end{tabular}
\end{table}

\section{Order of calls}
\begin{enumerate}
\item \cmd{CModel} constructor creates an empty \cmd{CGroundwaterModel}. The \cmd{.rvp} parser fills aquifer classes and
profiles; the \cmd{.rvh} parser stores \cmd{AQUIFER\_PROFILE} on each HRU; the \cmd{.rvt} parser stores well-rate and
boundary-head series; the \cmd{.rvc} parser stores initial or hotstart heads; \cmd{ParseGWFile} reads the \cmd{.rvg}.
\item \cmd{CModel::CalculateInitialWaterStorage} (called once, and again for each ensemble member) calls
\cmd{CGroundwaterModel::\allowbreak Initialize}: \cmd{CheckProcesses}, \cmd{BuildGeometry} (grid generation by \cmd{CGWGrid} and HRU--cell overlaps, with cache), \cmd{BuildVertical},
\cmd{BuildLinks}, \cmd{BuildRivers}, \cmd{BuildStresses} (wells, monitoring wells, GHB/CHD, surface stores, reservoirs),
\cmd{WriteMF6Files}, \cmd{ConnectEngine}, \cmd{OpenOutputs}, \cmd{WriteSummary}, \cmd{RunSteadyState}.
\item Every step, \cmd{MassEnergyBalance} (\cmd{Solvers.cpp}) calls \cmd{Exchange} after the lateral processes, then
\cmd{ApplyRiverExchange} before routing, and \cmd{TakeRiverLossFromInflow} inside the routing loop.
\cmd{IncrementCumulInput} adds \cmd{GetStepReturnVolume}; \cmd{UpdateDiagnostics} asks \cmd{GetObservationWellHead}.
\item \cmd{WriteMajorOutput} calls \cmd{WriteHotstart}; the destructor closes outputs and finalizes MODFLOW.
\end{enumerate}

\section{MODFLOW 6 memory used}
All are read or written through \cmd{get\_value\_ptr} after \cmd{prepare\_time\_step} (pointers are refreshed every step).
\begin{table}[!htbp]\centering
\caption{MODFLOW 6 memory used.}
\begin{tabular}{P{0.2\linewidth}P{0.388\linewidth}P{0.34\linewidth}}\toprule\hdr
Package (name) & Arrays & Use\\\midrule
GWF & \cmd{X} & heads\\
STO & \cmd{STRGSS}, \cmd{STRGSY} & storage flows\\
SLN\_1 & \cmd{MXITER} & iteration limit\\
RCH (\cmd{RCH\_RAVEN}) & \cmd{RECHARGE}, \cmd{HCOF}, \cmd{RHS} & recharge\\
DRN (\cmd{DRN\_RAVEN}) & \cmd{HCOF}, \cmd{RHS} (aux \cmd{DDRN} in file) & seepage\\
RIV (\cmd{RIV\_RAVEN}) & \cmd{STAGE}, \cmd{COND}, \cmd{RBOT}, \cmd{HCOF}, \cmd{RHS} & rivers\\
EVT (\cmd{EVT\_RAVEN}) & \cmd{RATE}, \cmd{HCOF}, \cmd{RHS} & water-table ET\\
WEL (\cmd{WEL\_RAVEN}) & \cmd{Q}, \cmd{HCOF}, \cmd{RHS} & wells\\
GHB (\cmd{GHB\_RAVEN}) & \cmd{BHEAD}, \cmd{COND}, \cmd{HCOF}, \cmd{RHS} & head boundaries\\
CHD (\cmd{CHD\_RAVEN}) & \cmd{HEAD}, \cmd{SIMVALS} & fixed heads\\
RIV (\cmd{SWX\_RAVEN}) & \cmd{STAGE}, \cmd{COND}, \cmd{RBOT}, \cmd{HCOF}, \cmd{RHS} & wetlands and lakes\\
WEL (\cmd{RES\_RAVEN}) & \cmd{Q}, \cmd{HCOF}, \cmd{RHS} & reservoir seepage\\\bottomrule
\end{tabular}
\end{table}
For MODFLOW versions before 6.5 the values are written to the columns of \cmd{BOUND} instead.

\section{Changes to existing Raven files}\label{sec:coreflix}
The changes are against Raven~4.15 (v593), and the branch \cmd{mf6-coupling} holds all of them. They fall into five
groups by what they change; the first is the coupling itself, the other four change Raven's behaviour also in runs
without groundwater:
\begin{enumerate}
\item \textbf{Groundwater coupling}: the new source files and the changes below that concern groundwater.
\item \textbf{Ensemble reset}: every ensemble member starts from the routing memory and
      reservoir states of the first run (\cmd{RestoreInitialDynamicState}) instead of inheriting the previous member's.
\item \textbf{Reservoir balance}: a drying reservoir books only the water available;
      reservoir stages and flows written to \cmd{solution.rvc} with 12 digits; reservoirs with absolute stages start at
      their crest when no initial condition is given; the reservoir constraint is initialised; a reservoir whose
      \cmd{:HRUID} lies in another subbasin stops with a message instead of crashing.
\item \textbf{Observation constituent}: \cmd{:ObservationData} and
      \cmd{:ObservationWeights} for \cmd{STREAM\_CONCENTRATION} read the constituent name before the line is re-read.
\item \textbf{Age-tracer hotstart}: the \cmd{:StartTime} offset of a hotstart reaches
      the age tracers. The masses keep the datum of the run that wrote the file, and the age clock and the entry time of
      precipitation are shifted by the offset. (Shifting the masses instead would make them negative, which the advection
      limiter cannot carry.) \cmd{solution.rvc} keeps \cmd{:StartTime} as the start of the run, as Raven~4.15 writes it,
      and adds a line \cmd{:AgeTracerDatum} with the datum when it is before that start (the start of the first run of a
      chain of hotstarts), so a chain keeps one datum. Reading a file, the offset is taken from \cmd{:AgeTracerDatum}, else
      from \cmd{:StartTime} when the file holds age-tracer masses, else it is 0 (a spin-up run without the tracer: the
      initial water has age 0). Runs without an age tracer, and cold starts, are unchanged.
\end{enumerate}
\begin{longtable}{P{0.32\linewidth}P{0.628\linewidth}}
\caption{Changes to existing Raven files (group in brackets).}\label{tab:corechanges}\\
\toprule\hdr
File & Change\\\midrule\endfirsthead
\caption[]{Changes to existing Raven files (continued).}\\
\toprule\hdr
File & Change\\\midrule\endhead
\cmd{Solvers.cpp} & groundwater hook after lateral processes; river exchange before routing; losses from reach inflow and outflow; with \cmd{:SoilZoneClasses} the HRU's process sequence once per column class (1)\\
\cmd{Model.cpp/.h} & always creates the groundwater object, deleted first so it can finish its outputs; returned water counted as input; \cmd{GW\_HEAD} diagnostics; flows of the column classes (\cmd{ClassBalance}, \cmd{SetFlowBalance}) (1); stored first routing state (2)\\
\cmd{ModelInitialize.cpp} & groundwater initialized after all inputs (1); routing memory and reservoir states of the first run snapshotted and restored before each member's .rvc is read (2); reservoir \cmd{:HRUID} check (3)\\
\cmd{StandardOutput.cpp} & groundwater hotstart blocks in \cmd{solution.rvc}; in a coupled run the HRU and basin states of \cmd{solution.rvc} with 17 significant digits (Raven's 5 decimals without groundwater); groundwater outputs closed at the end of the run (1); \cmd{:AgeTracerDatum} line in \cmd{solution.rvc} when the age tracers' datum is before the run's start; \cmd{:StartTime} unchanged (5)\\
\cmd{AgeTracers.cpp/.h}, \cmd{Advection.cpp} & the hotstart offset shifts the age clock, not the masses; precipitation enters at the shifted time; the datum kept for \cmd{solution.rvc} (\cmd{SetDatum}, \cmd{GetDatum}, \cmd{GetTimeOffset}) (5)\\
\cmd{ParseInput.cpp} & \cmd{:GroundwaterModel}; \cmd{:rvg\_Filename} relative to the \cmd{.rvi}; obsolete USG commands (\cmd{:Drain}, \cmd{:Recharge RECHARGE\_FLUX/RECHARGE\_DATA}) stop with clear errors; \cmd{:Transport} stops in a coupled run (the coupling exchanges water only, no constituent mass) (1)\\
\cmd{ParsePropertyFile.cpp} & \cmd{:AquiferClasses}, \cmd{:AquiferProfiles}, \cmd{:AquiferProfile\allowbreak Parameters} (1)\\
\cmd{ParseHRUFile.cpp}, \cmd{HydroUnits.h} & \cmd{AQUIFER\_PROFILE} stored on each HRU (1)\\
\cmd{ProcessGroup.h}, \cmd{SoilWaterMovers.h} & accessors the coupling reads (\cmd{GetSubProcess}, \cmd{GetPercolationType}) (1)\\
\cmd{ParseTimeSeriesFile.cpp} & \cmd{:WellFlowRate} (m$^3$/d, m$^3$/s or L/s), \cmd{:BoundaryHead} (m); other units stop (1); constituent name captured before parsing (4)\\
\cmd{ParseInitial\allowbreak ConditionFile.cpp} & \cmd{:InitialGWHeads} and the groundwater hotstart blocks (\cmd{:GWHeads}, \cmd{:GWStartHeads}, \cmd{:GWRechargeStore}, \cmd{:GWReservoirSeepage}, \cmd{:GWRiverLossCarry}, \cmd{:GWRiverTake}, \cmd{:GWSoilZoneColumns}, \cmd{:GWVadoseColumns}, \cmd{:GWVadoseClassColumns}, \cmd{:GWVadoseColumnTops}, \cmd{:GWSoilZoneLag}), each value checked (not a number as a whole token, not finite, or a negative water amount, stops: a damaged file; \cmd{:InitialGWHeads} read strictly, as the \cmd{.rvg}'s numbers), ignored with a warning in a run without \cmd{:GroundwaterModel}; the capacity check of the initial soil water skipped for the seepage target and the \cmd{:SurfaceWaterExchange} stores of a coupled model (they can hold more than their capacity, as in a continuous run) (1); ensemble members restored (2); default reservoir stage (3); \cmd{:AgeTracerDatum} read, age-tracer masses detected in the state table, and the offset passed to the age tracers before their initial conditions are set (5)\\
\cmd{SubBasin.cpp/.h} & channel depth at a given flow (1); last lateral inflow accessor (2)\\
\cmd{Reservoir.h/.cpp} & accessors for the seepage constant and local head (1); dynamic-state get/set (2); default initial stage, drying, \cmd{solution.rvc} precision (12 decimals, or the stream's 17 significant digits in a coupled run), constraint initialised (3)\\
\cmd{RavenInclude.h}, \cmd{CommonFunctions.cpp} & references to the MODFLOW-USG coupling removed (1); \cmd{RC\_UNSET} (3)\\
\cmd{CMakeLists.txt}, \cmd{src/Makefile}, \cmd{src/Raven.vcxproj(.filters)} & new files; link \cmd{-ldl} (CMake: \cmd{CMAKE\_DL\_LIBS}); CMake: Release by default for single-configuration generators, option \cmd{RAVEN\_MF6} with \cmd{cmake/RavenMF6*.cmake} (the MODFLOW~6 library: build, download or given path; install with its runtime) (1)\\
\cmd{README.md}, \cmd{.gitignore}, \cmd{.github/\allowbreak workflows/\allowbreak testing.yml} & the \cmd{RAVEN\_MF6} option described; \cmd{cmake/RavenMF6*} tracked; the existing test workflow configures with \cmd{-DRAVEN\_MF6=OFF} (and \cmd{.github/workflows/raven-mf6.yml} builds and tests with the library) (1)\\
\cmd{Drain.cpp}, \cmd{GWRecharge.cpp}, \cmd{GWRiverConnection.*}, \cmd{GWSWProcessABC.cpp}, \cmd{GWSWProcesses.h}, \cmd{MFUSGpp.h}, \cmd{RiverReach.*} & removed (1): the MODFLOW-USG coupling defined the same \cmd{CGroundwaterModel} class and only built with \cmd{\_MODFLOW\_USG\_} and the MF-USG library; it remains in the history of the Raven repository (v593)\\\bottomrule
\end{longtable}

\section{Where things live}
The coupling is split as the Raven team asked: Raven's own repository holds only the source code, and everything else
lives in a separate repository.
\begin{itemize}
\item \textbf{RavenHydroFramework} (\cmd{src/}): the coupling's source files and the small changes to existing Raven files
      (Section~\ref{sec:coreflix}). The branch \cmd{mf6-coupling} of
      \url{https://github.com/rarabzad/RavenHydroFramework} holds them all.
\item \textbf{RavenHydroFramework-extra}: this manual, the user guide, the Liard examples, the tests
      (\cmd{tests/run\_all.py}) and \cmd{tools/get\_mf6.py}, at\\ \url{https://github.com/rarabzad/RavenHydroFramework-extra}
      (branch \cmd{main}).
\end{itemize}
Clone the two side by side in one folder; the tests and \cmd{get\_mf6.py} then find Raven without any setting:
\begin{lstlisting}
git clone -b mf6-coupling https://github.com/rarabzad/RavenHydroFramework.git
git clone https://github.com/rarabzad/RavenHydroFramework-extra.git
\end{lstlisting}
Without network access, the release package's \cmd{source/} folder holds both repositories with their whole history as
git bundles (\cmd{tools/package\_sources.py}); the two \cmd{git clone} commands of its README (section ``Restore the git
history'') make the same two checkouts. \cmd{tools/make\_handover.py} writes the kit from which the two repositories are
published, with the steps to push and check them (\cmd{HANDOVER.md}). \cmd{tools/check\_manual\_pdf.py} builds this manual
from \cmd{docs/manual} in a scratch folder and checks that \cmd{docs/manual/main.pdf} is that build (the same text and
pages; \cmd{--write} replaces a PDF that is not); \cmd{package\_sources.py} runs it on the revision it packages.

\section{Building}\label{sec:building}
Raven is built with CMake (\cmd{cmake -S . -B build}, \cmd{cmake --build build --config Release}; with a
single-configuration generator such as Unix Makefiles or Ninja, which ignores \cmd{--config}, the build type defaults to
\cmd{Release}, \cmd{-DCMAKE\_BUILD\_TYPE=...} overrides it), with \cmd{make} in
\cmd{src/} (links \cmd{-ldl}; for NetCDF head output add \cmd{-Dnetcdf} to \cmd{CXXFLAGS} and \cmd{-lnetcdf} to
\cmd{LDLIBS}, as the Makefile's commented lines show), or on Windows with \cmd{src/Raven.sln}. Raven loads the MODFLOW~6
library at run time; with CMake the build can also provide that library and put it next to the executable, chosen by
the option \cmd{RAVEN\_MF6} (file \cmd{cmake/RavenMF6.cmake}):
\begin{table}[!htbp]\centering
\caption{The CMake option \cmd{RAVEN\_MF6}.}\label{tab:ravenmf6}
\begin{tabular}{P{0.13\linewidth}P{0.815\linewidth}}\toprule\hdr
Value & What the build does\\\midrule
\cmd{SOURCE} & builds \cmd{libmf6} from the MODFLOW~6 source with MODFLOW's own meson build, during \cmd{cmake --build}
(about one to three minutes on two cores). Needs a Fortran compiler (gfortran or ifx), meson~1.7 or newer, ninja and git.
The source is a shallow git clone of the release tag \cmd{RAVEN\_MF6\_VERSION} (6.8.1), checked against its pinned
commit (\cmd{5b2ebd1a95720d5c3732dc482e05fbf559e6ef41}); or a source folder (\cmd{RAVEN\_MF6\_SOURCE\_DIR}); or an
archive (\cmd{RAVEN\_MF6\_SOURCE\_URL}, also \cmd{file://}; a release zip works, its \cmd{src/} is the tag's, and a pinned
one is checked). The last two need no git.\\
\cmd{PREBUILT} & downloads the official release zip for the platform at configure time (\cmd{mf6.8.1\_linux.zip},
\cmd{mf6.8.1\_macarm.zip}, \cmd{mf6.8.1\_win64.zip}), checks it against the SHA-256 pinned in Raven and keeps only the
library. Offline: \cmd{-DRAVEN\_MF6\_ZIP=/path/to/mf6.8.1\_linux.zip}.\\
\cmd{PATH} & uses a library you have: \cmd{-DRAVEN\_MF6\_PATH=/path/to/libmf6.so} (a symbolic link or a renamed file such
as \cmd{libmf6\_663.so} or \cmd{libmf6.so.6.8.1} is fine).\\
\cmd{OFF} & Raven alone, with the same compile and link commands as without this option; \cmd{cmake --install}
installs Raven alone (\cmd{bin/Raven}), and the configure and the install say so. A coupled model then needs a library
at run time (Section~\ref{sec:getmf6}).\\\bottomrule
\end{tabular}
\end{table}

\paragraph{The default.} Without \cmd{RAVEN\_MF6} the choice is made again at every configure (the cache keeps only a value
you give): an offline option names its mode (\cmd{RAVEN\_MF6\_PATH} $\to$ \cmd{PATH}, \cmd{RAVEN\_MF6\_ZIP} $\to$
\cmd{PREBUILT}, \cmd{RAVEN\_MF6\_SOURCE\_DIR} or \cmd{RAVEN\_MF6\_SOURCE\_URL} $\to$ \cmd{SOURCE}); otherwise \cmd{SOURCE}
when a Fortran compiler, meson~$\geq$~1.7, ninja and git are all found \emph{and} the MODFLOW~6 git repository answers
at configure time (\cmd{git ls-remote} of the release tag, at most 60~s; not needed when the clone of an earlier configure is
already at the pinned commit); otherwise \cmd{PREBUILT} when MODFLOW releases a library for the platform; otherwise
\cmd{OFF} with a warning that names what is missing (the source tools, or, when they are installed, the network;
\cmd{-DRAVEN\_MF6=OFF} silences it). Every download the default needs
is thus made or tested at configure time: a default that cannot be completed (offline: the repository does not answer, or
the release zip cannot be downloaded) falls back with a warning (\cmd{SOURCE} $\to$ \cmd{PREBUILT} $\to$ \cmd{OFF}), so
neither a plain \cmd{cmake -S . -B build} nor the build after it stops because of MODFLOW; a mode you give (or one named
by an offline option) that cannot be completed stops with a message naming the alternatives (an explicit \cmd{SOURCE}
clones during the build and stops there if the clone fails).
\paragraph{The MODFLOW version.} A library older than 6.5.0 (\cmd{MF6\_MIN\_VERSION} in \cmd{src/MF6Engine.h}, the floor
of Raven's run-time guard) stops the configure when its version is known: \cmd{RAVEN\_MF6\_VERSION} for a clone or
download, the \cmd{meson.build} of a zip or source folder, the name of a source archive; another version than the tested
6.8.1 gives a warning. A \cmd{PATH} library's version is known only when Raven loads it (the run-time guard).
\paragraph{Installing.} \cmd{cmake --install} (and \cmd{cpack}) put Raven, \cmd{libmf6} and the Fortran runtime the library
needs into one \cmd{bin} folder (\cmd{cmake/RavenMF6Bundle.cmake}): on Linux the runtime libraries \cmd{ldd} lists, and
since Raven loads \cmd{libmf6} by its path, the installed \cmd{libmf6} must have the run path \cmd{\$ORIGIN} to load
them: a library built by this CMake from source has it, and one without it (for example a \cmd{PATH} library built
elsewhere) gets it from \cmd{patchelf} when that program is found; otherwise the install warns that the library loads
the system's runtime instead of the copies beside it; on macOS
every \cmd{libgfortran}, \cmd{libquadmath} and \cmd{libgcc\_s} the library needs directly or through another of them, with
\cmd{@rpath} references resolved through the \cmd{LC\_RPATH} entries, every reference rewritten to \cmd{@loader\_path} and the
edited files signed again (\cmd{codesign -f -s -}); on Windows the runtime DLLs the library imports and those they import
(\cmd{objdump}, else \cmd{dumpbin}). A \cmd{libmf6.dll} given with a folder (\cmd{:MF6Library}, \cmd{RAVEN\_MF6\_LIB}) is
loaded with its own folder searched first for the DLLs it needs. \cmd{cmake/RavenMF6InstallCheck.py <bin>} runs the
installed Raven on a tiny coupled model and checks that it used the library beside it and that every Fortran runtime
library it loaded lies in that folder (the CI job runs it after moving the install). An option that belongs to another mode is reported as
ignored.

\paragraph{Platforms without a release library.} MODFLOW~6.8.1 releases no library for Linux other than x86\_64 (e.g.\
aarch64), for Windows on ARM64 (the target architecture is taken from \cmd{-A ARM64}, else the compiler, else the
processor, and named in the message), or for Intel Macs (no Intel Mac zip from MODFLOW~6.8 on). There the default is
\cmd{SOURCE} when the tools are present, otherwise \cmd{OFF}; use \cmd{SOURCE} or \cmd{PATH}.

\paragraph{Checksums and offline use.} The release zips of 6.8.1 are pinned by SHA-256; only these are called ``official
release'' in the configure output. Another zip given by \cmd{RAVEN\_MF6\_ZIP}, or the download of another version, is used
with a warning (``not checked'') unless you give \cmd{-DRAVEN\_MF6\_ZIP\_SHA256=...}, which is then checked and named as
your checksum. A local zip is identified by its own SHA-256: a replaced zip (or another one with the same name) is
unpacked again, and the expected SHA-256 is compared with the file's at every configure, so a wrong checksum given later
stops the configure. A damaged or truncated zip is refused with the expected and the actual SHA-256.
Other settings: \cmd{RAVEN\_MF6\_VERSION} (release for \cmd{SOURCE} and \cmd{PREBUILT}; only 6.8.1 is pinned),
\cmd{RAVEN\_MESON}, \cmd{RAVEN\_MF6\_FORTRAN}, \cmd{RAVEN\_MF6\_JOBS}, \cmd{RAVEN\_MF6\_RELEASE\_URL},
\cmd{RAVEN\_MF6\_GIT\_URL}.

\paragraph{The Fortran compiler.} meson builds with \cmd{RAVEN\_MF6\_FORTRAN} if given, else with \cmd{FC} (the
environment variable meson itself reads), else with the compiler CMake's own check finds, else with gfortran or ifx on
the \cmd{PATH}. \cmd{FC} is read directly because CMake's check ignores it with the Visual Studio generators; the
compiler is named in the configure output and in \cmd{libmf6\_version.txt}. On Windows, gfortran from MSYS2 or
\cmd{fortran-lang/setup-fortran} works with the Visual Studio and the Ninja generators (the continuous-integration
workflow \cmd{.github/workflows/raven-mf6.yml} does this; it has not been run on real Windows and macOS runners).

\paragraph{What ends up where.} In every mode except \cmd{OFF} the library is copied next to the Raven executable under
its standard name (\cmd{libmf6.so}, \cmd{libmf6.dylib}, \cmd{libmf6.dll}), with \cmd{libmf6\_version.txt} beside it
(first line \cmd{MODFLOW 6.8.1} as written by \cmd{get\_mf6.py}, then its source and SHA-256). The version is the
release's, or that in the source's \cmd{meson.build} for source folders, archives and renamed zips (a tag archive named
\cmd{v6.7.2.tar.gz} gives 6.7.2 before it is built), and ``unknown (given library)'' for \cmd{PATH}. Raven's run path is
\cmd{\$ORIGIN} (macOS \cmd{@loader\_path}) plus the folders of linked libraries outside the system folders (netCDF from
conda or spack), in the build tree as in the install. \cmd{cmake --install build} puts Raven, the library,
\cmd{libmf6\_version.txt} and the Fortran runtime the library needs (a gfortran-built library: \cmd{libgfortran.so.5} on
Linux; official libraries need none) in one \cmd{bin} folder, also for staged installs with \cmd{DESTDIR}; the folder
runs wherever it is moved. \cmd{cpack} packs it as \cmd{Raven-4.15-mf6-6.8.1-linux-x86\_64.tar.gz} (about 5~MB with a
source-built library; ZIP on Windows; no MODFLOW version in the name when it is unknown).

\paragraph{Version guard.} After loading the library, Raven asks its version (\cmd{get\_version}). A library older
than 6.5.0 stops the run (``the MODFLOW 6 library \ldots\ is version \ldots; the coupling needs MODFLOW 6 6.5.0 or newer
(tested with 6.8.1)''); 6.5.0 is a chosen floor (6.4.4 passed test P1; the fallbacks for older memory layouts are
kept). Any other version than 6.8.1, a library that does not report its version, or a version that cannot be read
(quoted in the message) gives one warning; outputs are otherwise unchanged. Suffixes such as \cmd{6.8.1-dev} or
\cmd{6.8.1+abc.dirty} (a source build inside a git checkout) count as 6.8.1.

\paragraph{Checks.} \cmd{tests/cmake\_checks.py} runs 22 offline configure checks of the above in about a minute, each
also shown to fail on the defect it guards against (\cmd{-{}-selftest}; Table~\ref{tab:vcmake}, Section~\ref{sec:vbuild});
\cmd{tests/mf6version\_check.py} checks the version guard with a stand-in library. A library built from source and the official one give the same results to
round-off (Section~\ref{sec:getmf6}).

\section{Getting MODFLOW 6}\label{sec:getmf6}
A Raven built with \cmd{RAVEN\_MF6} other than \cmd{OFF} has the library next to it and needs nothing more. Otherwise
the MODFLOW~6 shared library (\cmd{libmf6.so}, \cmd{libmf6.dll} or \cmd{libmf6.dylib}, about 15--55~MB) is not part of
either repository. One command in RavenHydroFramework-extra downloads it:
\begin{lstlisting}
python tools/get_mf6.py                   # MODFLOW 6.8.1, the tested version, for this computer
python tools/get_mf6.py --latest          # the newest release instead (or --version 6.x.y)
python tools/get_mf6.py --zip mf6.8.1_win64.zip   # from a release zip downloaded by hand (offline)
python tools/get_mf6.py --raven PATH      # PATH: the Raven checkout, or the folder of Raven.exe
\end{lstlisting}
The script (standard Python only) picks the release zip for the computer from
\url{https://github.com/MODFLOW-ORG/modflow6/releases} (\cmd{mf6.8.1\_linux.zip}, \cmd{mf6.8.1\_win64.zip},
\cmd{mf6.8.1\_macarm.zip}), keeps only the library, and records its version, source and SHA-256 in
\cmd{libmf6\_version.txt} beside it. It puts the library in \cmd{lib/mf6/} of the Raven checkout beside
RavenHydroFramework-extra (or the one named by \cmd{--raven} or \cmd{RAVEN\_SRC}), where Raven finds it. It does nothing
when the library is already there (\cmd{--force} replaces it). It serves no library where MODFLOW releases none: Linux
other than x86\_64 and Windows on ARM64 (it stops and points to \cmd{cmake -DRAVEN\_MF6=SOURCE} or \cmd{PATH}; on Windows
on ARM64 a win64 zip given with \cmd{--zip} is used, with a note, for an x64 Raven running under emulation), and Intel
Macs from MODFLOW~6.8 on (\cmd{--version 6.6.3}, the last tested release with an Intel Mac zip, or \cmd{SOURCE}). The
default is the tested version rather than the newest release, so that a new MODFLOW release cannot change a model's
results unnoticed; with MODFLOW~6.6, for example, the Newton option behaved differently (Section~\ref{sec:mf6versions}).

Raven looks for the library in this order: \cmd{:MF6Library} in the \cmd{.rvg} file; the environment variable
\cmd{RAVEN\_MF6\_LIB}; next to the Raven executable; \cmd{lib/mf6/} beside it or one to three folders up (so
\cmd{src/Raven.exe}, \cmd{build/Raven}, \cmd{build/Release/Raven.exe} and Visual Studio's \cmd{src/x64/Release/Raven.exe}
of a Raven checkout all find the checkout's \cmd{lib/mf6/}); finally the platform's name on the system library path.
\cmd{GWModelSummary.txt} names the library used and its version; when none loads, the message lists the places searched.

\paragraph{Built from source or downloaded.} A library built here from the 6.8.1 source (gfortran~13, meson's release
build) and the official one give the same results to round-off: in a one-year run of the Liard example the heads differ by at most $9\times10^{-11}$~m, and \cmd{Hydrographs.csv}, \cmd{WatershedStorage.csv} and \cmd{GWHRUState.csv} are identical. Building the library takes about one minute on two cores here; the installed folder with a source-built library (about 5~MB packed) ran this example after being moved, with no setting.

\section{Adding a new exchange}
\begin{enumerate}
\item Choose the MODFLOW package (a new package name, e.g.\ \cmd{XYZ\_RAVEN}) and decide which Raven quantity drives it.
\item Parse its inputs in \cmd{ParseGWFile.cpp} (or the \cmd{.rvp}/\cmd{.rvt} parsers) and store them in \cmd{CGroundwaterModel}.
\item Select its cells in \cmd{BuildStresses}; write the package file and its line in \cmd{gwf.nam} in \cmd{WriteMF6Files}.
\item Fetch its arrays in \cmd{RefreshPointers} (named arrays first, \cmd{BOUND} as fallback).
\item Set its values in \cmd{Exchange} (or \cmd{SetBoundaryArrays}) after \cmd{prepare\_time\_step}.
\item After the solve, compute its flow as $\mathrm{HCOF}\cdot h-\mathrm{RHS}$; apply the matching volume on the Raven side
and add it to \cmd{GetStepReturnVolume} if it changes Raven's storage.
\item Add the flow to the balance error $\varepsilon$, the gross-flow scale, the \cmd{GWBudget.csv} header and line, and the
non-convergence cap if it is an inflow.
\item Add a test, extend \cmd{tools\_audit.py} and the list-file cross-check.
\end{enumerate}
\section{Adding a grid type}
A grid type only has to produce cells: add a generator to \cmd{CGWGrid} that fills the cell polygons (counterclockwise, in the
grid's local frame) and, if it knows them exactly, the neighbours; \cmd{Finish} computes areas, centres, bounding boxes,
spatial bins and welded vertices, and finds neighbours from shared edges otherwise. Choose the MODFLOW format in
\cmd{Initialize} (DISU works for any cells) and the XT3D default in \cmd{UseXT3D}. Everything else --- overlaps, rivers,
boundaries, wells, outputs, restarts --- works through the cell polygons and needs no change. Add the new type to the grid unit
test (\cmd{tests/gridtest.cpp}), the Thiem test (\cmd{tests/thiem/}) and the randomized tests.

\section{Testing procedure}
\begin{enumerate}
\item Uncoupled regression: run a model without groundwater with the new build and compare \cmd{Hydrographs.csv} and
\cmd{Watershed\allowbreak Storage.csv} byte for byte with a build without your change (\cmd{RAVEN\_EXE\_REF}).
\item Build with \cmd{-Wall -Wextra -Wshadow} and run cppcheck on the new files.
\item Run the feature tests and \cmd{tools\_audit.py} on each output folder.
\item Compare with MODFLOW's list-file budget.
\item Build with \cmd{-fsanitize=address,undefined} and run the feature tests.
\end{enumerate}

\section{Limitations and known issues}\label{ch:limits}
\paragraph{Current limitations.}
\begin{itemize}
\item Refinement of a regular grid spans whole rows and columns; use a quadtree or HRU mesh for local refinement.
\item XT3D (optional) makes each iteration more expensive and can slow the Newton method badly where cells dry and rewet; it is off by default.
\item Imported grid files must hold one single-ring polygon per cell (MODFLOW cells have no holes).
\item Input files are GeoJSON, ESRI ASCII grids and GeoTIFF files (Section~\ref{sec:geotiff}: single band, uncompressed,
LZW, DEFLATE or PackBits, north-up), in longitude/latitude or in one projection given by \cmd{:InputCRS}
(Section~\ref{sec:inputcrs}: Transverse Mercator, Albers and Lambert conformal conic on WGS84 or GRS80; other projections
and other raster encodings must be converted first).
\item DEM, bedrock, reservoir stages, boundary heads and well screens must use one vertical datum; the only automatic check
is the stop for a coupled reservoir whose stages lie below the aquifer under its lake.
\item By default layer $\ell$ connects sideways only to layer $\ell$ of neighbouring columns; \cmd{:LayerConnection ELEVATION}
connects the cells whose elevations overlap (Section~\ref{sec:layerconn}); an existing model keeps its own connections.
\item River stage uses Raven's channel depth at the start of the step; reservoir seepage reaches the aquifer one step
later (both exactly accounted).
\item A layer that starts completely dry can cause isolated non-converged steps; Raven warns at start-up when a whole layer
starts dry (Section~\ref{ch:layers}). Start from a steady state or realistic heads.
\item MODFLOW stops the program on invalid input (Fortran \cmd{stop}); Raven copies the error report of
\cmd{mf6/mfsim.lst} into its message and \cmd{Raven\_errors.txt}.
\item Between cells, soil water moves only through the aquifer (\cmd{:SoilZone PHYSICAL}); seepage faces smaller than a cell
need smaller cells.
\item Existing models (Chapter~\ref{ch:existing}, Section~\ref{sec:extlimits}): no DISU grids; no ATS6, several solutions or
a non-Gregorian calendar; restarts and later starts refuse time-array series (TAS6), TVK6/TVS6, simulations with several
models, array or advanced packages with several period entries before the start, and \cmd{TS6} files whose rows come
from another file (\cmd{TS6} time series are shifted to Raven's start); TDIS and water-mover period data given as \cmd{INTERNAL} are refused (\cmd{OPEN/CLOSE} files are read); the stream
package keeps the model's stage and has no conductance (supply) limit; undefined \cmd{LENGTH\_UNITS}/\cmd{TIME\_UNITS} are
taken as metres and days (set them in DIS and TDIS); monitoring wells and \cmd{GWHeads.nc} need the cell outlines
(\cmd{:MF6CRS} or \cmd{:MF6GridFile}, not \cmd{:OverlapWeights}). Water-table ET, \cmd{:SurfaceWaterExchange} and
\cmd{:ReservoirExchange} work as with a generated model. In a model without \cmd{NEWTON} a top cell that dries becomes
inactive: MODFLOW moves recharge and evapotranspiration to the active cell below, but not Raven's seepage drain or a
reservoir's seepage, which stop on that column until the cell rewets (the reservoir's seepage is returned to Raven and
booked); the water table in the layer below can then rise above the ground without seeping out. Use \cmd{NEWTON} where top
cells dry.
\end{itemize}
\paragraph{Known issues in Raven itself.} The following belong to Raven's own code, behave the same without groundwater, and
are left to the Raven developers:
\begin{itemize}
\item Raven's channel-storage bookkeeping (\cmd{:ChannelStorage} in \cmd{solution.rvc}) can be negative in some subbasins
(17 subbasins of the Liard model without groundwater).
\item Reservoir routine: on the day after a reservoir has been drawn empty by a minimum flow, the outflow still averages in
the release rate at the end of the previous day, although the lake is empty; this makes up most of such a reservoir's balance
error (Section~\ref{sec:behaviour}, Reservoirs that run dry).
\item Reservoir routine: the stage iteration stops once its change falls below 0.1~mm. On a step whose stage ends within a
millimetre of the weir crest, where the weir's outflow curve has its kink, the remaining imbalance is up to about the lake
area times 0.1~mm (a few m$^3$ on a 0.04~km$^2$ lake), shown as the reservoir's \emph{MB error} in
\cmd{ReservoirMassBalance.csv} and in Raven's balance. Base Raven gives the same without groundwater (3 to 7~m$^3$ on such a
lake whose stage falls slowly through the crest, with a fixed \cmd{:SeepageParameters} level); with \cmd{:ReservoirExchange}
it appears where the aquifer holds a lake near its crest (Section~\ref{sec:vcombinations}).
\item Hotstart of a model whose reservoir spills: Raven's balance error (\emph{MB Error} of \cmd{WatershedStorage.csv}) of a
run started from \cmd{solution.rvc} is offset, from its first step on and by a constant, by half the water that crossed the
reservoir in the step before the restart: half its outflow and, where reservoir seepage returned to the subbasin or a gain
MODFLOW did not supply flows on into the reservoir (Section~\ref{ch:reservoirs}), half that inflow. Raven's initial storage
at a hotstart leaves this half step out; base Raven 4.15 does the same without groundwater. On a 0.16~km$^2$ test watershed
whose lake spills 2\,364~m$^3$ in the step before the restart, the offset is $+7.39$~mm (1\,182~m$^3$) without groundwater,
in base Raven and with the coupling alike, and $+14.92$~mm coupled (half of 4\,776~m$^3$); a lake spilling 16\,100~m$^3$/d gives about
50~mm; returned or unsupplied reservoir seepage in transit gives 0.1 to 0.6~mm. The restarted run's states, flows and stages
are the continuous run's (exactly in a coupled run, to the digits of Raven's \cmd{solution.rvc} without groundwater); only
this column is offset, and with it the Raven balance that \cmd{tools\_audit.py} reports (\emph{Raven MB error max}).
\item Raven's balance error (\emph{MB Error} of \cmd{WatershedStorage.csv}) of the first step is not zero in some models and
stays at that value afterwards: 0.431~mm in the Liard example with daily steps and 0.216~mm with half-day steps, in base
Raven~4.15 without groundwater and in a coupled run alike (Section~\ref{sec:accraven}).
\item Without \cmd{:GroundwaterModel}, Raven empties the \cmd{GROUNDWATER} store at the start of every step without booking
its water as an output (base Raven books it only for model types other than a surface-water model), so a model that sends
water to \cmd{GROUNDWATER} shows a balance error that falls by the store's content every step: $-7.48$~mm after the first
year of the Liard example without groundwater. The flows are not affected.
\item In ensemble mode, \cmd{WatershedStorage.csv} shows channel and rivulet storages that differ from a single run, also
without groundwater; the hydrographs, reservoir stages and groundwater outputs of each member equal the single run's
(Section~\ref{sec:behaviour}).
\item valgrind reports uninitialised reads in base Raven when it reads the \cmd{.rvp} file: in \cmd{RVPParameterWarning}
(called from \cmd{ParseClassPropertiesFile}; three contexts) and in \cmd{CVegetationClass::RecalculateRootParams} (two
contexts; 23 errors in the short run of test I10). The coupling's own code and the hotstart code show none
(Section~\ref{sec:codeaudit}).
\end{itemize}
\paragraph{Not available.}
Groundwater transport (MODFLOW~6 GWT) linked to Raven constituents; particle tracking (PRT) during the run (offline tracking
on the coupled flows is possible: \cmd{:WriteGWBudgetFile} and \cmd{tools/make\_prt.py}, Section~\ref{sec:budgetfile}; PRT
needs DIS or DISV grids); wells driven by Raven water demands (irrigation). Frozen ground acts on recharge only through Raven's own
processes (the coupling receives whatever they send to \cmd{GROUNDWATER}).
